% CVPR 2026 Paper Template; see https://github.com/cvpr-org/author-kit

\documentclass[10pt,twocolumn,letterpaper]{article}

%%%%%%%%% PAPER TYPE  - PLEASE UPDATE FOR FINAL VERSION
% \usepackage{cvpr}              % To produce the CAMERA-READY version
\usepackage[review]{cvpr}      % To produce the REVIEW version
% \usepackage[pagenumbers]{cvpr} % To force page numbers, e.g. for an arXiv version

\usepackage[utf8]{inputenc} % allow utf-8 input
\usepackage[T1]{fontenc}    % use 8-bit T1 fonts
\usepackage{url}            % simple URL typesetting
\usepackage{booktabs}       % professional-quality tables
\usepackage{amsfonts}       % blackboard math symbols
\usepackage{nicefrac}       % compact symbols for 1/2, etc.
\usepackage{microtype}      % microtypography
\usepackage{xcolor}         % colors
\usepackage{comment}
\usepackage{graphicx}
\usepackage{multirow}
\usepackage{enumitem}
\usepackage{tabularx}
\usepackage{pifont}

%%
%% Inline annotations; for predefined colors, refer to "dvipsnames" in the xcolor package:
%% https://tinyurl.com/overleaf-colors
%%
\newcommand{\todo}[1]{\textbf{\color{red}[TODO: #1]}}
\newcommand{\sam}[1]{\textbf{\color{blue}[Sam: #1]}}
\newcommand{\jake}[1]{\textbf{\color{green}[Jake: #1]}}
\newcommand{\gpt}[1]{\textbf{\color{orange}[GPT: #1]}}

\usepackage[detect-all,separate-uncertainty = true]{siunitx}
\sisetup{
    output-exponent-marker=\ensuremath{\mathrm{e}},
    group-separator= {,},
    group-minimum-digits = 4,
    list-final-separator={, },
    mode={math},
    retain-explicit-plus
}

\makeatletter
\renewcommand{\paragraph}{%
  \@startsection{paragraph}{4}%
  {\z@}{1.5ex \@plus 1ex \@minus .2ex}{-0.5em}%
  {\normalfont\normalsize\bfseries}%
}
\makeatother

\setlist{nosep}

%%
%% globally adjusts space between figure and caption
%%
\setlength{\abovecaptionskip}{.5em}


%%
%% Commonly used math definitions
%%
% \DeclareMathOperator*{\argmin}{arg\,min}
% \DeclareMathOperator*{\argmax}{arg\,max}

%% Tigthen underline
\usepackage{soul}
\setuldepth{foobar}

% It is strongly recommended to use hyperref, especially for the review version.
% hyperref with option pagebackref eases the reviewers' job.
% Please disable hyperref *only* if you encounter grave issues, e.g. with the file validation for the camera-ready version.
\definecolor{cvprblue}{rgb}{0.21,0.49,0.74}
\usepackage[pagebackref,breaklinks,colorlinks,allcolors=cvprblue]{hyperref}
\usepackage[capitalize]{cleveref}

%%%%%%%%% PAPER ID  - PLEASE UPDATE
\def\paperID{*****} % *** Enter the Paper ID here
\def\confName{CVPR}
\def\confYear{2026}

\definecolor{green}{RGB}{25, 169, 116}
\definecolor{red}{RGB}{255, 65, 54}
\newcommand{\cmark}{\textcolor{green}{\ding{51}}}
\newcommand{\xmark}{\textcolor{red}{\ding{55}}}
% For fishvista
\definecolor{blue}{RGB}{0, 18, 115}
\definecolor{cyan}{RGB}{10, 147, 150}
\definecolor{sea}{RGB}{148, 210, 189}

%%%%%%%%% TITLE - PLEASE UPDATE
\title{Sparse Autoencoders for Data-Driven Morphological Trait Discovery}

\author{Samuel Stevens\\
The Ohio State University\\
\texttt{stevens.994@osu.edu}\\
\and
Jacob Beattie\\
The Ohio State University\\
\texttt{beattie.74@osu.edu}\\
\and
% Neil Rosser\\
% University of Miami\\
% % Wildlife?\\
% \texttt{rosser?}\\
% \and
% Daniel Rubenstien\\
% Princeton University\\
% % EEB?\\
% \texttt{dir@princeton.edu}\\
% \and
Tanya Berger-Wolf\\
The Ohio State University\\
% Computer Science \& Engineering\\
\texttt{berger-wolf.1@osu.edu}\\
\and
Yu Su\\
The Ohio State University\\
% Computer Science \& Engineering\\
\texttt{su.809@osu.edu}
}

\bibliographystyle{plainnat}


\begin{document}

\maketitle

\begin{abstract}
\sam{Need a sentence talking about the large scientific datasets first.}
Large, weakly labeled scientific datasets have driven the development of foundation models whose internal representations encode structure (patterns, co-occurrences and statistical regularities) beyond task objectives.
Current methods for extracting this structure from this models require pre-specifying target concepts; these approaches excel at confirmation but cannot perform open-ended exploration to discover patterns that lack established names or annotations.
We address this limitation by applying sparse autoencoders (SAEs) to decompose foundation model representations into interpretable features without concept supervision.
In controlled rediscovery studies, the learned features align with semantic concepts on a standard segmentation benchmark and surpass strong label-free alternatives on concept-alignment metrics.
Applied to ecological imagery, the same procedure surfaces fine-grained anatomical structure without access to segmentation or part labels.
While our experiments focus on vision with an ecology case study, the method is domain-agnostic and directly portable to models in other sciences (e.g., proteins, genomics, weather).
The results provide a practical route to turn foundation-model representations into reusable measurements for discovery.
\end{abstract}

\section{Introduction}

Scientific data now span hundreds of petabytes across domains, far outpacing what any lab can annotate or manually inspect.\footnote{For example, the NIH Sequence Read Archive exposes roughly 47 PB of public sequencing data, the European Nucleotide Archive provides roughly 63 PB of immediately downloadable data, NASA Earthdata (EOSDIS) hosts more than 120 PB of open Earth science data, and a single reanalysis, ERA5, is on the order of 5 PB \citep{SRA2024,ENA2024,NASA_EOSDIS,ERA5}.}
Large neural networks trained on these datasets cover proteins (e.g., AlphaFold 3 \citep{abramson2024alphafold3}), weather (e.g., GraphCast, GenCast \citep{lam2023graphcast,price2025gencast}), genomics (e.g., scGPT, ESM3 \citep{cui2024scgpt,hayes2025esm3}), and vision (e.g., Scale-MAE, BioCLIP 2 \citep{reed2023scale,gu2025bioclip2}).
These models achieve strong predictive performance across a variety of tasks by learning compressed representations of the underlying data distribution.
% They learn to organize data by recurring patterns and factors \citep{bommasoni2021opportunities}; in vision, self-supervised encoders demonstrably learn object- and part-level structure that transfers to dense prediction tasks \citep{caron2021dinov1,oquab2024dinov2}.

\begin{figure*}[t]
    \centering
    \includegraphics[width=\linewidth]{figures/hook.png}
    % \includegraphics[width=\linewidth]{figures/hook.pdf}
    \caption{
    }\label{fig:hook}
\end{figure*}

However, current usage of such foundation models in science is predominantly task-centric: models are used as frozen feature extractors or fine-tuned for specific predictions (species classification \citep{vanhorn2021inat21}, cellular responses to perturbations \citep{roohani2025vcc}, 3-D protein structure \citep{haas2018continuous,pereira2021high}).
The learned, compressed representations are high-dimensional embeddings optimized for task performance.
While these representations enable accurate predictions, they do not directly provide reusable scientific measurements or expose the structure the model learned for systematic investigation.
This limits foundation models to \emph{confirmatory} applications where the target concepts are pre-specified (see \cref{tab:related-work}).

\paragraph{Discovery vs. Confirmation}
Scientific discovery often involves finding patterns that lack established names or annotations.
\sam{Need a citation.} \jake{A proposal for this paragraph: Scientific discovery often involves the creation of hypotheses that rely on known concepts and patterns. (no need to cite here, I'd argue it's common knowledge. However, researchers often want to formalize previously unknown concepts or uncover unknown relationships. (the next sentence stays)}
A morphological feature that distinguishes subspecies but has not been formally catalogued, a correlation between climate variables that no prior study examined, or a conserved sequence motif in an understudied protein family: these cannot be found by methods that require specifying target concepts in advance.
\sam{Need some citations. There are a lot of claims here that are not trivially true.} \jake{Recently, AI has been used to uncover new drugs \citep{sadybekov_computational_2023}, gene associations with rare diseases \citep{cipriani_rare_2025}, and new molecules \citep{bushuiev_self-supervised_2025}. (Really, our sell isnt that this problem is unsolved. Lots of people already do scientific discovery with ai. The gap we're trying to fill is visual/morphological concept discovery.)}

Foundation models trained on millions of samples have processed more data than any human could learn from.
If these models have learned semantic patterns not captured by existing annotations, interpretability methods offer a route to systematically surface this learned structure.
The patterns are encoded in the model's representations: model interpretability offers a path towards unlocking model-learned discoveries.

However, current interpretability approaches require specifying concepts, limiting us to confirmatory work:
\begin{itemize}
\item{Supervised probing tests for known concepts \citep{alain2016understanding}.}
\item{Concept activation vectors (CAV) measure sensitivity to user-defined concepts \citep{kim2018tcav}.}
\item{Saliency maps and attention visualization explain individual predictions for predefined tasks \citep{selvaraju2017gradcam,simonyan2013deep}.}
\end{itemize}
These methods excel at \textit{confirmation}, but cannot perform open-ended exploration to generate hypotheses about what patterns exist in the learned representations.
\todo{Discuss neural network visualization where you do gradient descent on the image.} \jake{Integrated gradients/gradients on images could probably be listed alongside saliency maps and attention visualization.}
When labels are scarce and salient factors are unknown (i.e., in most scientific use cases \todo{cite something here}), the inability to extract unnamed patterns from foundation models represents a significant limitation.
If foundation models have learned meaningful patterns from data that exceed what humans can manually analyze, we need methods to surface those patterns without pre-specification.

\begin{table*}[t]
    \centering
    \small
    \renewcommand{\arraystretch}{1.2}
    \setlength{\tabcolsep}{4pt}
    \caption{Comparison of methods for extracting knowledge from foundation models. Most approaches require researchers to specify concepts or tasks in advance, limiting them to confirmatory analysis. SAEs enable open-ended discovery by surfacing the model's learned feature vocabulary without supervision.}
    \label{tab:related-work}
    \begin{tabular}{lllll}
    \toprule
    \textbf{Method} & \textbf{Input} & \textbf{Output} & \textbf{Example} & \textbf{Discovery?} \\
    \midrule
    Task Heads & Labeled task dataset & Task predictions & Species classification & \xmark{} Needs task labels \\
    Linear Probes & Labeled concept $\mathcal{A}$ samples & Does model know concept $\mathcal{A}$? & ``Does DINOv3 know eyes?'' & \xmark{} Pre-specified concepts \\
    Concept Vec. & Labeled concept $\mathcal{A}$ samples & Degree of concept $\mathcal{A}$ & ``Find `striped' neuron'' & \xmark{} Needs concept samples \\
    Feature Viz. & \\
    Prototypes & Labeled task dataset & Interpretable predictions\\
    SAEs & Unlabeled data & Interpretable features & Discover new traits & \cmark{} Finds unknown patterns \\
    \bottomrule
    \end{tabular}
\end{table*}

\paragraph{Our approach.}
We address this gap by applying sparse autoencoders (SAEs) to decompose foundation model representations into interpretable features without concept supervision \citep{cunningham2024sparse,gao2025scaling}.
SAEs learn an overcomplete dictionary that reconstructs model activations through sparse linear combinations, with sparsity encouraging individual features to capture coherent semantic units.
Unlike supervised methods, SAEs require no concept labels during training.
Each learned feature provides a per-example activation score together with its decoding direction and exemplar evidence, enabling systematic investigation of what patterns the model represents.
In short, SAEs are a promising approach for discovery from foundation models \citep{peng2025unknown}.

\paragraph{Evidence.}
We train SAEs on patch-level activations from DINOv3 \citep{simeoni2025dinov3}, a self-supervised vision transformer (ViT, \cite{dosovitskiy2020vit}).
We evaluate if SAEs can recover fine-grained semantic structure without access to labels in two settings: (1) ADE20K scene segmentation \citep{zhou2017ade20k}, containing 150 object classes, and (2) ecological images, where we test recovery of anatomical body parts from the FishVista dataset \citep{mehrab2024fishvista,fishvistadataset}.
We find that SAEs reliably extract semantic concepts that were never explicitly labeled in DINOv3's self-supervised training objective.
We compare against decomposition baselines ($k$-means clustering, PCA, non-negative matrix factorization) and find that SAEs achieve substantially higher concept alignment, with a Matryoshka SAE variant \citep{bussmann2025matryoshka} reaching \num{70.4}\% higher class coverage than standard SAEs on ADE20K (see \cref{sec:ade20k,tab:baselines} for details).

Our results indicate that unsupervised sparse decomposition can systematically recover semantic structure from foundation model representations.
On standard vision benchmarks, SAEs rediscover segmentation classes and spatial relationships without label access during training.
On scientific images, they surface fine-grained anatomical features, providing measurable quantities that can be evaluated against biological annotations.
We analyze when this approach succeeds and when it fails, documenting sensitivity to architectural choices, layer selection, and dataset characteristics.
While we validate on computer vision, the method applies to any domain with token-level foundation model representations: proteins (residues), genomics (genes), weather (grid cells), and other structured data \citep{abramson2024alphafold3,lam2023graphcast,cui2024scgpt}.
\sam{We need 1-2 more ``pie-in-the-sky'' sentences to remind readers of our vision.}

\section{Related Work}

Theory and evidence suggest that large-pretrained models internalize meaningful structure \citep{ha2018world,caron2021dinov1,lecun2022jepa,huh2024platonic} but prior work develops highly-specific probes rather than discovery engines.

\paragraph{Foundation Models.}
Domain-specific foundation models capture rich structure in their modalities and achieve state-of-the-art performance on their focal tasks, as demonstrated in proteins \citep{abramson2024alphafold3,hayes2025esm3}, ecology \citep{stevens2024bioclip}, weather \citep{lam2023graphcast,price2025gencast}, and single-cell genomics \citep{cui2024scgpt}.
In both general and scientific domains, large-scale training leads to emergent abilities beyond the pre-training objective \citep{wei2022emergent,gu2025bioclip2}.
Despite this, most scientific models are benchmarked as \textbf{task-specific predictors} rather than sources of novel concepts.

\paragraph{Interpretability Methods.}
A large body of work explains predictions without exposing a global factorization of representations.
Saliency and attribution methods (e.g., gradients, Integrated Gradients, Grad-CAM) highlight input regions for \emph{individual} decisions and have known validity pitfalls \citep{simonyan2014deep,sundararajan2017axiomatic,selvaraju2017gradcam,adebayo2018sanity}.
Linear probes test separability of \emph{pre-specified} targets \citep{alain2017probes,hewitt2019control}; concept methods such as TCAV/ACE require exemplar sets and so test investigator-defined ideas \citep{kim2018tcav,ghorbani2019ace}.
Prototype models (ProtoPNet, ProtoTree, ProtoPFormer) provide case-based, class-tied explanations but presuppose a labeled concept space \citep{chen2019protopnet,nauta2021prototree,protopformer2022}.
Activation maximization and “feature visualization” (DeepDream-style) optimize inputs for chosen units to aid intuition, but they do not yield dataset-level, per-example measurements \citep{mahendran2016visualize,nguyen2016dgn,olah2017feature}.
Unsupervised dense-discovery methods (LOST, TokenCut, STEGO, MaskCut) reveal structure via clustering or segmentation without labels \citep{lost2021,ji2021tokencut,stego2022,maskcut2022}, yet typically return masks or clusters rather than a shared library of features.

Instance-level methods such as saliency maps and GradCAM \citep{simonyan2014deep,selvaraju2017gradcam} explain individual predictions by highlighting input regions, but are sensitive to design choices \citep{adebayo2018sanity,sturmfels2020visualizing}.
Feature visualization methods generate synthetic inputs that maximally activate specific neurons \citep{olah2017feature,mordvintsev2015deepdream}, revealing what patterns networks respond to, but \textbf{require pre-selecting which neurons to visualize} and produce synthetic rather than real-data measurements.
Concept-based methods like TCAV and ACE \citep{kim2018tcav,ghorbani2019ace} measure sensitivity to user-defined concepts, while network dissection \citep{bau2017network} aligns neurons with predefined ontologies.
Linear probing tests whether specific concepts are linearly decodable from representations \citep{alain2016understanding}, though this conflates probe capacity with representational content \citep{hewitt2019probes,belinkov2022probing}.
Prototype-based methods like ProtoPNet \citep{chen2019protopnet} learn interpretable classifiers by matching to prototypical training examples, but are trained end-to-end for specific tasks and \textbf{require labeled data for predefined classes}.
These approaches excel at explaining predictions or testing for known concepts (confirmation) but cannot be used to discover patterns without prior specification.

\paragraph{Sparse Autoencoders \& Decomposition.}
Classical dictionary learning and sparse coding demonstrated that sparse, parts-based factors can emerge from natural data \citep{olshausen1996nature,aharon2006ksvd,lee1999nmf}.
Recent work applies sparse autoencoders (SAEs) to model activations, showing that sparse decompositions can partially ``unmix'' superposed features and yield monosemantic directions \citep{elhage2022superposition,bricken2023monosemantic,gao2024scaling} in language and vision models.
SAEs address the superposition hypothesis \citep{elhage2022superposition}—that networks represent more features than they have neurons by assigning features to directions in activation space rather than individual neurons.
Gao et al.\ \citep{gao2024scaling} demonstrate that k-sparse autoencoders scale effectively and improve reconstruction-sparsity tradeoffs.
Concurrent work explores SAEs for discovering unknown concepts rather than acting on known ones \citep{peng2025use} and applies dictionary learning to microscopy foundation models for biological concept extraction \citep{zimmermann2024scientific}.
However, systematic evaluation of SAEs for semantic concept recovery on standard benchmarks, comparison to classical decomposition baselines, and application to scientific discovery with ground-truth validation remains limited.

\paragraph{Discovery.}
We use \emph{discovery} to mean surfacing previously unnamed factors (morphological features distinguishing subspecies, correlations between climate variables, or conserved sequence motifs in understudied proteins) as opposed to \emph{confirmation} of pre-specified targets.
Closest in intent to our goal are (i) work advocating SAEs for uncovering \emph{unknown} concepts rather than acting on \emph{known} ones \citep{peng2025unknown} and (ii) dictionary-learning approaches that extract biological concepts from microscopy foundation models \citep{microscopy_dictionary_learning_2025}.
We differ by targeting dense token representations in vision and validating both on a standard semantic benchmark and an ecology case study.
We also differ from unsupervised dense segmentation methods \citep{lost2021,ji2021tokencut,stego2022,maskcut2022} by returning a reusable \emph{measurement library} (global features with per-example activations) rather than per-image masks or clusters.

Methods for systematic, unsupervised discovery from foundation model representations remain underdeveloped.
Our work addresses this gap by demonstrating that sparse decomposition of foundation model representations can systematically surface semantic structure.



\section{Methodology}

We ask the question: ``Can SAEs discover unlabeled concepts in pre-trained models?''
Because discovery demands finding unnamed structure rather than confirming known categories, our method avoids concept supervision and treats labels only as holdout probes.
Our goal is to recover a dictionary whose latents serve as candidate concepts and other interpretable features that can be measured and compared across images.
We outline the ViT feature extraction, the SAE objective, and the evaluation probes that operationalize this goal.

\paragraph{Foundation Model.} We use the publicly available DINOv3 ViT-L/16 checkpoint \citep{simeoni2025dinov3}.
Images are resized to $256\times256$ with bicubic interpolation, yielding a $16\times16$ grid of patches per image.
We extract patch-token activations ([CLS] discarded). 
We expect the later half of ViT layers to contain the best features \citep{bolya2025perceptionencoder}; in all experiments, we sweep over \num{6} evenly-spaced layers: for \num{24}-layer ViT-L models, we use layers \num{14}, \num{16}, \num{18}, \num{20}, \num{22}, and \num{24}.\footnote{For \num{12}-layer ViT-S and ViT-B models, we use the last \num{6} layers: \num{7}, \num{8}, \num{9}, \num{10}, \num{11}, and \num{12}.}

\paragraph{Sparse Autoencoders (SAEs).} Our default dictionary learner is a Matryoshka SAE trained independently for each layer. 
We train \num{16384}-latent SAEs for 100M tokens and sweep a logarithmic grid of sparsity coefficients $\lambda$ and learning rates.
We summarize each trained model by its reconstruction quality (normalized MSE) and sparsity (L0).
% normalized_mse (float): reconstruction sum squared SAE error / reconstruction sum squared mean baseline error
Normalized MSE (NMSE) is MSE normalized by a mean baseline MSE:
\begin{equation}\label{eq:nmse}
\text{NMSE} = \frac{\sum_i || x_i - \hat x_i ||_2^2}{\sum_i || x_i - \bar x ||_2^2},
\end{equation}
where $\hat x_i$ is the predicted reconstruction and $\bar x$ is the mean of all $x_i$.
\begin{figure}[t]
    \centering
    \small
    \includegraphics[width=\linewidth]{figures/dinov3_vitl16_in1k_sae.pdf}
    \caption{We compare $k$-means, PCA, vanilla SAEs and Matryoshka SAEs along the reconstruction-sparsity tradeoff for learning to decompose ViT patch activations. We use final layer activations from DINOv3 ViT-L/16 on ImageNet-1K; we fit all methods on the training split and measure normalized MSE (see \cref{eq:nmse}) and L0 on the validation split. For $k$-means, we ``reconstruct'' every patch with its nearest centroid; thus, L0 is always $1$. For PCA, we sweep the number of non-zero components used during inference. For both SAE variants, we sweep $\lambda$ and show the Pareto frontier of reconstruction-sparsity. \textbf{Takeaway:} Reconstruction-sparsity does not indicate an optimal method.
    \sam{PCA results are not done yet.}
    }
    \label{fig:mse-l0}
\end{figure}


\paragraph{Matryoshka SAEs.}
To avoid feature splitting \citep{bricken2023monosemanticity,chanin2024absorption,leask2025sparse}, we use Matryoshka SAEs \cite{bussmann2025matryoshka}, which learn a single SAE by training multiple nested SAEs simultaneously.
Each training step samples random prefixes of SAE latents and tries to reconstruct the inputs using a limited number of latents.
We sample ten prefixes $m_1, m_2, ... m_10 \in [1, 16384]$, so the first $m_1$ latents must reconstruct the input as best they can, the first $m_2$ latents must do the same, etc.
This forces early latents to capture more general features.

\begin{figure*}[t]
    \small
    \centering
    \begin{subfigure}[b]{0.32\textwidth}
        \centering
        \includegraphics[width=0.49\textwidth]{figures/kmeans-ade20k-sky-1.png}
        \hfill
        \includegraphics[width=0.49\textwidth]{figures/kmeans-ade20k-sky-2.png}
        \includegraphics[width=0.49\textwidth]{figures/kmeans-ade20k-sky-3.png}
        \hfill
        \includegraphics[width=0.49\textwidth]{figures/kmeans-ade20k-sky-4.png}
        \includegraphics[width=0.49\textwidth]{figures/kmeans-ade20k-sky-5.png}
        \hfill
        \includegraphics[width=0.49\textwidth]{figures/kmeans-ade20k-sky-6.png}
        \caption{$k$-means}
        \label{fig:ade20k-examples-kmeans}
    \end{subfigure}
    \hfill
    \begin{subfigure}[b]{0.32\textwidth}
        \centering
        \includegraphics[width=0.49\textwidth]{figures/pca-ade20k-sky-1.png}
        \hfill
        \includegraphics[width=0.49\textwidth]{figures/pca-ade20k-sky-2.png}
        \includegraphics[width=0.49\textwidth]{figures/pca-ade20k-sky-3.png}
        \hfill
        \includegraphics[width=0.49\textwidth]{figures/pca-ade20k-sky-4.png}
        \includegraphics[width=0.49\textwidth]{figures/pca-ade20k-sky-5.png}
        \hfill
        \includegraphics[width=0.49\textwidth]{figures/pca-ade20k-sky-6.png}
        \caption{PCA}
        \label{fig:ade20k-examples-pca}
    \end{subfigure}
    \hfill
    \begin{subfigure}[b]{0.32\textwidth}
        \centering
        \includegraphics[width=0.49\textwidth]{figures/sae-ade20k-toilet-1.png}
        \hfill
        \includegraphics[width=0.49\textwidth]{figures/sae-ade20k-toilet-2.png}
        \includegraphics[width=0.49\textwidth]{figures/sae-ade20k-toilet-3.png}
        \hfill
        \includegraphics[width=0.49\textwidth]{figures/sae-ade20k-toilet-4.png}
        \includegraphics[width=0.49\textwidth]{figures/sae-ade20k-toilet-5.png}
        \hfill
        \includegraphics[width=0.49\textwidth]{figures/sae-ade20k-toilet-6.png}
        \caption{Matryoshka SAE}
        \label{fig:ade20k-examples-sae}
    \end{subfigure}
    \caption{
    \textbf{Takeaway}: \todo{Fill the takeaway out.}
    }
    \label{fig:ade20k-examples}
\end{figure*}


\paragraph{Evaluation.}
\sam{Need one sentence explaining why we need evaluation.}
We evaluate discovery capability by \textit{re}-discovering ADE20K \citep{zhou2017ade20k} semantic classes via 1D logistic probes, using the same protocol as \cite{gao2025scaling}.
% Following the protocol, a patch is labeled positive for class c if the corresponding segmentation mask covers at least 25\% of the patch area. 
For each candidate SAE and each ADE20K class, we fit ridge-regularized logistic regressions independently for every latent against the binary patch labels on ADE20K train and select the best latent per class by train loss. 
The primary metric is Probe R (pseudo-$R^2$ on cross-entropy), defined as the mean relative reduction in binary cross-entropy versus a prevalence-only baseline; we also report mAP, Purity@$k$ (precision over the top-$k$ latent-scored patches for $k=16$), and Coverage@$\tau$ (fraction of classes whose best AP $\geq \tau$ for $\tau=0.3$). 
Dictionary-learning metrics (NMSE, L0) are reported alongside downstream metrics.




\section{Experiments}

\sam{Need an introduction sentence.}

\subsection{Re-Discovering Semantic Segmentations}\label{sec:ade20k}

\sam{Need an introduction sentence.}

We train sparse autoencoders (SAEs) on a neutral, unlabeled corpus. 
Concretely, we extract activations on ImageNet-1K (train/val) for every \{backbone, layer\} pair. 
Per-layer feature means and standard deviations are computed on ImageNet-1K and used to standardize activations at training and evaluation time. 
Separately, we extract and store ADE20K activations (train/val) for probing and reporting.

\begin{table}[t]
    \centering
    \small
    \setlength{\tabcolsep}{3pt}
    \caption{SAEs for re-discovering ADE20K's semantic segmentation classes from DINOv3 ViT-L/16's pre-trained representations of ImageNet-1K.
    For each method, we choose the best layer from DINOv3 based on probe loss on ADE20K's training split.
    % Results for all layers are included in \cref{app:ade20k-all}.
    Probe $R$ is explained variance, as described in \cref{sec:ade20k}.
    Purity@$k$ is precision over the top-$k$ highest-scoring patches for a feature, where $k=16$.
    Coverage@$\tau$ is the fraction of ADE20K classes with best AP $\geq\tau$, where $\tau= 0.3$.
    \textbf{Takeaway:} Matryoshka SAEs demonstrate meaningful improvement over baseline methods on all metrics, and outperform vanilla SAEs on all downstream metrics despite worse dictionary learning metrics (NMSE and sparsity).}
    \label{tab:ade20k}
    \begin{tabular}{lrrrrrr}
    \toprule
    & \multicolumn{2}{c}{Dict. $\downarrow$} & \multicolumn{4}{c}{Downstream $\uparrow$} \\
    \cmidrule(lr){2-3} \cmidrule(lr){4-7}
    Method & NMSE & L0 & Probe $R$ & mAP & Purity@$k$ & Cov@$\tau$ \\
    \midrule
    $k$-Means & 0.619 & 1 & 0.001 & 0.031 & 0.795 & 0.026 \\  % myy5btgw, layer 24/24
    PCA & 0.687 & 64 & 0.236 & 0.065 & 0.591 & 0.079 \\ % a1x1laxm, 24/24
    % NMF \\
    SAE        & 0.460 & 64.1 & 0.425 & 0.313 & & 0.417 \\  % 0pdum8cq, layer 24/24
    Matryoshka & 0.702 &   6.5 & 0.402 & 0.319 & 0.808 & 0.450 \\  % mccrm7u8, layer 24/24
    \bottomrule
    \end{tabular}
\end{table}

\paragraph{Baseline Methods.}
We compare SAEs to two standard unsupervised machine learning methods: $k$-means and principal component analysis (PCA).
We fit $k$-means ($k$=\num{16384}) and PCA (\num{1024} components).
$k$-means uses one-hot codes with L$0$=$1$; PCA keeps the top $n$ coefficients ($n \in \{1,4,16,64,256,1024\}$).
We report normalized reconstruction MSE on ImageNet-1K (validation split).
% See \cref{app:baselines} for more details.

\paragraph{Results.}
We compare all methods on the reconstruction-sparsity tradeoff in \cref{fig:mse-l0}, report concept rediscovery metrics in \cref{tab:ade20k} and show qualitative examples in of rediscovered concepts in \cref{fig:ade20k-examples}.
While both SAE variants do not improve upon classical baselines as measured by normalized MSE and sparsity, they \textit{do} lead to better concept discovery as measured by mAP, Purity@$k$ and Coverage@$\tau$.
\todo{Anecdotally, PCA and Matryoshka SAE latents are easier to \textit{browse} because they are \textit{ordered}; each ``latent'' from $k$-means and vanilla SAEs are equally weighted, while PCA and Matryoshka latents are ordered by ``generality''.}




\subsection{Can SAEs recover domain-specific concepts?}
\label{sec:fishvista}

Beyond the general domain semantic concepts present in ADE20K, we want to use SAEs to enable scientific discovery.
Again, we evaluate discovery capability via \textit{re}-discovery: can vision foundation models, in combination with SAEs, re-discover the anatomy of diverse fish species? 
We use FishVista \citep{mehrab2024fishvista,fishvistadataset} and the 10 annotated body-part traits to answer this question (see \cref{fig:fishvista} for dataset examples).
We train SAEs on \num{56.3}K fish images and evaluate SAE discovery on the 4.6K (disjoint) images with body part segmentations.


\begin{figure*}[t]
    \small
    \centering
    \begin{subfigure}[b]{0.32\textwidth}
        \centering
        \includegraphics[width=0.49\textwidth]{figures/dinov3-fishvista-head-sae-1.png}
        \hfill
        \includegraphics[width=0.49\textwidth]{figures/dinov3-fishvista-head-seg-1.png}
        \includegraphics[width=0.49\textwidth]{figures/dinov3-fishvista-head-sae-2.png}
        \hfill
        \includegraphics[width=0.49\textwidth]{figures/dinov3-fishvista-head-seg-2.png}
        \includegraphics[width=0.49\textwidth]{figures/dinov3-fishvista-head-sae-3.png}
        \hfill
        \includegraphics[width=0.49\textwidth]{figures/dinov3-fishvista-head-seg-3.png}
        \caption{\textbf{\textcolor{blue}{``Head'' (\texttt{DINOv3-16K/9})}}}
        \label{fig:fishvista-head}
    \end{subfigure}
    \hfill
        \begin{subfigure}[b]{0.32\textwidth}
        \centering
        \includegraphics[width=0.49\textwidth]{figures/dinov3-fishvista-dorsal-sae-1.png}
        \hfill
        \includegraphics[width=0.49\textwidth]{figures/dinov3-fishvista-dorsal-seg-1.png}
        \includegraphics[width=0.49\textwidth]{figures/dinov3-fishvista-dorsal-sae-2.png}
        \hfill
        \includegraphics[width=0.49\textwidth]{figures/dinov3-fishvista-dorsal-seg-2.png}
        \includegraphics[width=0.49\textwidth]{figures/dinov3-fishvista-dorsal-sae-3.png}
        \hfill
        \includegraphics[width=0.49\textwidth]{figures/dinov3-fishvista-dorsal-seg-3.png}
        \caption{\textbf{\textcolor{sea}{``Dorsal Fin'' (\texttt{DINOv3-16K/13})}}}
        \label{fig:fishvista-dorsal}
    \end{subfigure}
    \hfill
        \begin{subfigure}[b]{0.32\textwidth}
        \centering
        \includegraphics[width=0.49\textwidth]{figures/dinov3-fishvista-eye-sae-3.png}
        \hfill
        \includegraphics[width=0.49\textwidth]{figures/dinov3-fishvista-eye-seg-3.png}
        \includegraphics[width=0.49\textwidth]{figures/dinov3-fishvista-eye-sae-4.png}
        \hfill
        \includegraphics[width=0.49\textwidth]{figures/dinov3-fishvista-eye-seg-4.png}
        \includegraphics[width=0.49\textwidth]{figures/dinov3-fishvista-eye-sae-5.png}
        \hfill
        \includegraphics[width=0.49\textwidth]{figures/dinov3-fishvista-eye-seg-5.png}
        \caption{\textbf{\textcolor{cyan}{``Eye'' (\texttt{DINOv3-16K/69})}}}
        \label{fig:fishvista-eye}
    \end{subfigure}
    \caption{
    Example SAE features (left) and patch-level segmentation masks (right) for \textbf{\textcolor{blue}{(a) ``Head''}}, \textbf{\textcolor{sea}{(b) ``Dorsal Fin''}} and \textbf{\textcolor{cyan}{(c) ``Eye''}} from DINOv3.
    Features were picked by minimizing cross entropy on binary classification for each body part, as described in \cref{sec:fishvista}.
    \textbf{Takeaway}: \todo{Fill the takeaway out.}
    }
    \label{fig:fishvista}
\end{figure*}

\paragraph{Results.}
We report results in \cref{tab:fishvista} and include qualitative examples of three body parts (heads, eyes, and dorsal fins) in \cref{fig:fishvista}.
Matryoshka SAEs again outperfom all methods on downstream metrics.
Qualitatively, a small fraction of images are reversed (fish head on the right, rather than the left); despite this, the SAE features are consistent across these variations.
These results indicate that strong foundation models, in combination with SAEs, can learn scientifically sound features.
Furthermore, the Matryoshka objective naturally orders latents from most to least general; we (anecdotally, but consistently) found earlier latents more interesting than later latents.
\sam{Can you plot the purity@k with respect to the latent? Matryoshka indicates an order; this would.}

\begin{table}[t]
    \centering
    \small
    \setlength{\tabcolsep}{3pt}
    \caption{SAEs for re-discovering FishVista's ten fish body parts from DINOv3 ViT-L/16's pre-trained representations of ImageNet-1K.
    For each method, we choose the last layer from DINOv3.
    % Results for all layers are included in \cref{app:fishvista-all}.
    Probe $R$ is explained variance, as described in \cref{sec:fishvista}.
    Purity@$k$ is precision over the top-$k$ highest-scoring patches for a feature, where $k=16$.
    Coverage@$\tau$ is the fraction of FishVista body parts with best AP $\geq\tau$, where $\tau= 0.3$.
    \textbf{Takeaway:} Matryoshka SAEs demonstrate meaningful improvement over baseline methods on all metrics, and outperform vanilla SAEs on all downstream metrics despite worse dictionary learning metrics (NMSE and sparsity).}
    \label{tab:fishvista}
    \begin{tabular}{lrrrrrr}
    \toprule
    & \multicolumn{2}{c}{Dict. $\downarrow$} & \multicolumn{4}{c}{Downstream $\uparrow$} \\
    \cmidrule(lr){2-3} \cmidrule(lr){4-7}
    Method & NMSE & L0 & Probe $R$ & mAP & Purity@$k$ & Cov@$\tau$ \\
    \midrule
    $k$-Means & \\
    PCA \\
    % NMF \\
    SAE        & & & & & & \\
    Matryoshka & 0.150 & 75.1 & 0.441 & 0.610 & & 0.9 \\
    \bottomrule
    \end{tabular}
\end{table}


% For size, I am doing ViT-S, ViT-B and ViT-L on IN1K/ADE20K.
% For width, I'm happy to do 1024 x {2, 4, 8, 16, 32} with ViT-L on IN1K/ADE20K.
% For resolution I'm using a ViT-S for storage reasons and am doing 256 patches/img (default), then 576, then 1024. All of them use a native resize that picks out a size that minimizes aspect ratio distortion. So the 256 patch images can be 16 x 16 patches, 8 x 32 patches, 32 x 8 patches, etc. The 576 can be 24 x 24 patches, 32 x 18 patches, 18 x 32 patches, 48 x 12 etc. Then the 1024 can be 32 x 32 patches, etc. 
% 
% Remember, each of these ablations need to measure two things: dictionary learning quality (for completeness) and downstream metrics, BECAUSE WE WANT TO ENABLE DISCOVERY. I think it's very important to emphasize the effects of these ablations on discovery-related metrics rather than NMSE/L0.



\subsection{Do larger models discover better?}
We assess model size effects by training SAEs on patch activations from DINOv3 ViT-S/16, ViT-B/16, and ViT-L/16 checkpoints.
The models vary in both total parameter count (22M, 86M and 303M, respectively) and the embedding dimension $d$ (\num{384}, \num{768}, and \num{1024}, respectively).
\sam{Because of the difference in embedding dimension, I think smaller $d$ should lead to sparse SAEs, because there's a lower-dimension orthogonal basis or something. @Jake it would be great to get your math background here.}
For each size, we train SAEs with 16K latents on ImageNet-1K activations and probe the SAEs for ADE20K features.

\paragraph{Results.}
In \cref{fig:ablations-size-mse} we see that larger models are in fact harder to reconstruct (note that the $y$-axis is \textit{raw} rather than normalized MSE), consistent with prior work in scaling SAEs for language models \citep{gao2025scaling}.
Despite more challenging dictionary learning, we see in \cref{tab:vit-size} that SAEs applied to larger models both recover more concepts (coverage) and are more precise (purity).
All three sizes of the DINOv3 ViTs are distilled from a single, larger teacher model \citep{simeoni2025dinov3}; we conjecture then that all three checkpoints have learned the same set of concepts.
\sam{I don't know what to make of this. Why is it easier to recover from larger ViTs?}
These results indicate that larger pre-trained models, in combination with SAEs, are more effective for scientific discovery.
% \todo{Motivated by this, we train SAEs on DINOv3 ViT-7B/16.}


\begin{figure}[t]
    \centering
    \small
    \includegraphics[width=\linewidth]{figures/dinov3_sizes_in1k_sae.pdf}
    \caption{The reconstruction-sparsity tradeoff for SAEs trained on different sized ViTs. Note that the $y$-axis is \textit{raw} rather than normalized MSE. 
    \textbf{Takeaway:} \todo{Fill in.}}
    \label{fig:ablations-size-mse}
\end{figure}

\begin{table}[t]
    \centering
    \small
    \setlength{\tabcolsep}{3pt}
    \caption{
        SAEs on different sizes of DINOv3 for re-discovering ADE20K's semantic segmentation classes.
        \textbf{Takeaway:} While larger ViTs are harder to sparsify, they lead to better concept (re)discovery.
    }
    \label{tab:vit-size}
    \begin{tabular}{lrrrrrr}
    \toprule
    \multirow{2.4}*{ViT} & \multicolumn{2}{c}{Dict. $\downarrow$} & \multicolumn{4}{c}{Downstream $\uparrow$} \\
    \cmidrule(lr){2-3} \cmidrule(lr){4-7}
     & NMSE & L0 & Probe $R$ & mAP & Purity@$k$ & Cov@$\tau$ \\
    \midrule
    ViT-S/16 & \textbf{0.293} & 95.0 & \textbf{0.433} & 0.203 & 0.742 & 0.245 \\  % gc6iqrf2, layer 12/12
    ViT-B/16 & 0.359 & 89.8 & 0.440 & 0.312 & \textbf{0.811} & 0.430 \\  % qoc1660r, layer 12/12
    ViT-L/16 & 0.702 & \textbf{6.5} & 0.402 & \textbf{0.319} & 0.808 & \textbf{0.450} \\  % mccrm7u8, layer 24/24
    \bottomrule
    \end{tabular}
\end{table}


\subsection{Which layer is best?}
Representations evolve across network depth: later layers in CNNs and ViTs tend to encode higher-level, more linearly separable semantics, which often improves transfer and dense alignment \citep{yosinski2014transfer,azizpour2014factors,caron2021dinov1}. 
However, recent results show that the final layer is not always optimal \citep{bolya2025perceptionencoder} .
Intermediate layers can yield less entangled, more spatially localized features depending on the downstream task. 
We therefore evaluate 6 layers in the second half of the transformer for the three evaluated checkpoints (DINOv3 ViT-S/16, ViT-B/16 and ViT-L/16).

\paragraph{Results.}
In \cref{fig:ablations-layers-mse}, we find that later layers are harder to reconstruct (note that the $y$-axis is \textit{raw} rather than normalized MSE), 
Priot work in SAEs for language models finds that later layers are also harder to reconstruct except for the final two layers \citep{gao2025scaling}.  % fig 27 in appendix
We conjecture that language model's linear token embedding head optimize for linear separability in the final layer, leading to easier reconstruction.
In contrast DINOv3 uses multiple, loss-specific MLPs and so there is no reason that the final layer should be easy to reconstruct.
\sam{Very clunky. @Jake what do you think of this?}
We also evaluate downstream metrics (mAP) for the best SAE at each layer for all three ViT checkpoints in \cref{fig:ablations-best-layer}.
Despite more challenging reconstruction, we find that later layers lead to better recovery for all three checkpoints.

\begin{figure}[t]
    \centering
    \small
    \includegraphics[width=\linewidth]{figures/dinov3_vitl_layers_in1k_sae.pdf}
    \caption{The reconstruction-sparsity tradeoff for different layers of DINOv3 ViT-L/16. Again, note that the $y$-axis is \textit{raw} rather than normalized MSE.
    \textbf{Takeaway:} \todo{Fill in.}}
    \label{fig:ablations-layers-mse}
\end{figure}


\begin{figure}[t]
    \centering
    \small
    \includegraphics[width=\linewidth]{figures/dinov3_in1k_ade20k_layers_map.pdf}
    \caption{Matyroshka SAE probe quality as measured by mean average precision (mAP) for DINOv3 ViT-S/16, ViT-B/16 and ViT-L/16 checkpoints over different transformer layers. 
    \textbf{Takeaway:} \todo{Fill in.}}
    \label{fig:ablations-best-layer}
\end{figure}




\begin{figure*}[t]
    \centering
    \includegraphics[width=\linewidth]{figures/dinov3_vits16_in1k_ade20k_probe1d_nmse_map.pdf}
    \includegraphics[width=\linewidth]{figures/dinov3_vitb16_in1k_ade20k_probe1d_nmse_map.pdf}
    \includegraphics[width=\linewidth]{figures/dinov3_vitl16_in1k_ade20k_probe1d_nmse_map.pdf}
    \caption{\textbf{Top to bottom:} Matyroshka SAE probe quality for DINOv3 ViT-S/16, ViT-B/16 and ViT-L/16 checkpoints for 6 layers in the second half. \textbf{Takeaway:} \todo{Fill in.}}
    \label{fig:vit-size}
\end{figure*}

% \paragraph{Do larger SAEs enable better discovery?}
% We sweep the SAE dictionary size using ViT-L/16 patch activations to study capacity versus sparsity.
% Concretely, we set the number of latents to $S=1024\times m$ with $m\in\{2,4,8,16,32\}$, i.e., $S\in\{2\mathrm{K},4\mathrm{K},8\mathrm{K},16\mathrm{K},32\mathrm{K}\}$.
% We keep the encoder fixed (DINOv3 ViT-L/16).

% In \cref{fig:sae-width}
% \begin{figure}[t]
%     \centering
%     \small
%     \includegraphics[width=\linewidth]{figures/sae-width.png}
%     \caption{Bigger SAEs are better.}
%     \label{fig:sae-width}
% \end{figure}

\section{Conclusion \& Discussion}








% \cite{tsutsumi2023morphovae} use a VAE \todo{cite that} on primate jaws and find (existing) morphological traits. 
% The advantage of the VAE is that you can generate images from the latent space to precisely visualize the what each dimension means.
% \todo{Read and see if/how they formalize the problem}

% \cite{manogaran2024you} doesn't push hard enough on discovering novelty.
% It develops a method, but doesn't check whether you learn new things.

% \cite{elhamod2023discovering} doesn't actually discover new traits, despite the title.

% \cite{nauta2021neural}: novel method for interpretable classification, still no focus on discovering NEW traits.

% % This article tries to use a different method for decomposing MLP activations besides SAEs. They do better on steering, but about the same on concept recovery. We are solely interested in concept recovery
% \cite{shafran2025decomposing}: ``We find that features learned with SNMF outperform widely used SAEs (Lieberum et al., 2024; He et al., 2024) and supervised features computed with difference-in-means (Marks and Tegmark, 2024; Rimsky et al., 2024; Turner et al., 2024) on concept steering, while performing comparably on concept detection.''


\clearpage

\bibliography{main}

\clearpage

\appendix

\section{Baseline Methods}\label{app:baselines}

\begin{figure*}[t]
    \centering
    \includegraphics[width=\linewidth]{figures/dinov3_vits16_in1k_ade20k_probe1d_logl0.pdf}
    \includegraphics[width=\linewidth]{figures/dinov3_vitb16_in1k_ade20k_probe1d_logl0.pdf}
    \includegraphics[width=\linewidth]{figures/dinov3_vitl16_in1k_ade20k_probe1d_logl0.pdf}
    \caption{\textbf{Top to bottom:} Matyroshka SAE probe quality for DINOv3 ViT-S/16, ViT-B/16 and ViT-L/16 checkpoints for 6 layers in the second half. \textbf{Takeaway:} \todo{Fill in.}}
    \label{fig:vit-size-logl0}
\end{figure*}

\begin{figure}[t]
    \centering
    \includegraphics[width=\linewidth]{figures/dinov3_in1k_ade20k_layers.pdf}
    \caption{\textbf{Top to bottom:} Matyroshka SAE probe quality for DINOv3 ViT-S/16, ViT-B/16 and ViT-L/16 checkpoints for 6 layers in the second half. \textbf{Takeaway:} \todo{Fill in.}}
    \label{fig:vit-layers-probe-r}
\end{figure}

\end{document}

